\section{Embedding Models for Retrieval}
\medskip

Classification predicts discrete labels. Retrieval ranks documents by relevance.

\bigskip
\subsection{Bi-Encoder Architecture}
\medskip

In a bi-encoder, queries and documents are encoded independently:
\begin{align*}
v_Q &= f_\theta(Q), \\
v_{\kappa_i} &= g_\theta(\kappa_i).
\end{align*}

Relevance is computed by inner product:
\[
\text{score}(Q, \kappa_i)
=
\langle v_Q, v_{\kappa_i} \rangle.
\]

The advantage is scalability. Document embeddings can be computed once and stored. At query time, only $v_Q$ must be computed. Retrieval becomes a nearest-neighbor search in vector space.

\bigskip
\subsection{Contrastive Fine-Tuning with InfoNCE}
\medskip

To optimize retrieval directly, we fine-tune using labeled pairs $\{(q_i, d_i)\}$.

For a batch of size $B$, we compute embeddings
\begin{align*}
u_i &= f_{\theta_Q}(q_i), \\
v_j &= g_{\theta_D}(d_j),
\end{align*}
and similarities
\[
S_{ij} = \langle u_i, v_j \rangle.
\]

The InfoNCE loss is
\[
\mathcal{L}_i
=
- \log
\frac{\exp(S_{ii})}
{\sum_{k=1}^{B} \exp(S_{ik})},
\quad
\mathcal{L}
=
\frac{1}{B}
\sum_{i=1}^B \mathcal{L}_i.
\]

This objective pulls matching pairs together and pushes mismatched pairs apart. Because negatives are drawn from the same batch, the model learns relative discrimination rather than absolute scoring.

\bigskip

